% =====================================================================
%  SNIPPET: headline results for the introduction. DRAFTED 2026-07-07;
%  EXTENDED same day per JS to seven items: theta_QCD rephrased with
%  his rationale (first fruit of the GHZ-quark identification; "a hard
%  problem ... becomes trivial in the language of entanglement"), and
%  the dimension-of-space item added in his own voice at glimpse
%  status, op:three-dimensions attached so the hedge is priced.
%  Ordering: the dimension item sits sixth, not last, so the list
%  still closes on the falsifiable neutrino bet (JS to veto).
%  Selection criterion: a number against a measurement, one line each,
%  no double-counting (W/Z ratio excluded as the same fact as
%  sin^2 theta_W; the 8/9 spectrum rides the wrapper; glueball excluded
%  at its own consistency status; structural matches left to the prose).
%  NOTE: this is the paper's second deliberate exception to the
%  no-lists rule (the neutrino table is the first) -- the side-by-side
%  is the argument. JS to veto if preferred as prose.
%  All \refs verified to exist. British English.
% =====================================================================

From the chain summarised above, with zero adjustable dimensionless parameters and the electroweak
vacuum value $v$ as the single dimensional input, the framework returns --- among a spectrum in
which eight of the nine charged-fermion masses land within about one per cent of observation ---
the following headline results.

\begin{itemize}
\item $\sin^{2}\theta_{W} = \tfrac14$ at tree level, running to $0.2312$ against the measured
      $0.23122$ (\S\ref{sec:boson-masses}).
\item $m_{H} = v/2$: $123.1$~GeV against the measured $125.25$~GeV, within $1.7\%$
      (\S\ref{sec:boson-masses}).
\item $y_{t} = 1$ exactly: $m_{t} = v/\sqrt{2} = 174.1$~GeV against the measured $172.5$~GeV,
      within $0.9\%$ (\S\ref{sec:boson-masses}).
\item The Koide charged-lepton relation with its single angle fixed at $\delta_{0} = \tfrac29$,
      matching to about a tenth of a per cent (\S\ref{sec:masses}).
\item $\theta_{\mathrm{QCD}} = 0$ exactly --- the first consequence to fall out of identifying
      quarks with the GHZ class: the colour sector's topology admits no CP-violating phase, so a
      hard problem of QCD becomes trivial when expressed in the language of entanglement, no axion
      required, against the bound $\lvert\theta\rvert < 10^{-10}$
      (\S\ref{sec:entanglement-classes}).
\item By defining spacelike separation in terms of non-computability, the framework glimpses a
      rationale for the three dimensions of space, although a proof is as yet elusive
      (Open Problem~\ref{op:three-dimensions}; \S\ref{sec:space-of-computations}).
\item The lightest neutrino exactly massless, the ordering normal, and
      $\sum m_{\nu} \approx 52$~meV: consistent with every current bound and decidable by JUNO
      within a few years (\S\ref{sec:predictions}).
\end{itemize}

Where a derived number disagrees with experiment, the paper says so in the same table as its
successes (\S\ref{sec:predictions}); the statuses --- proved, structural, owed, open --- are
attached where each result is derived and are never silently upgraded.
